\subsection{Explicit progress on the three-pointed torus}

\begin{result}{D.28}{Proved}{The exact word problem, flag index and four stabilizer families for $(1,3)$}
Write $a_1,a_2,a_3$ for three mutually disjoint parallel
nonseparating curves, with one marked point in each intervening
annulus, and let $b$ meet each once.  Label $p_i$ as the point
opposite $a_i$ in the standard star picture.  There is the
explicit presentation
\[
 \Gamma_{1,3}=
 \left\langle a_1,a_2,a_3,b\ \middle|\
 [a_i,a_j]=1,\ a_i b a_i=b a_i b,\
 \Delta=(a_1a_2a_3b)^3=1\right\rangle.                       \tag{D.22}
\]
The maps send the four letters to exactly these Dehn twists.
The executed classical Garside algorithm in
\path{code/d_artin_d4.py} solves its word problem by computing
in $A(D_4)$ and retaining the proper greedy factors after
removing the central $\Delta$ exponent.

For $\{i,j,k\}=\{1,2,3\}$, put
\[
 d_i=(a_jba_k)^4,\qquad q_i=a_jba_k a_jba_j.
\]
Then $d_i$ is the separating twist enclosing $p_j,p_k$,
and the following are exact stabilizers with literal
embeddings into (D.22):
\[
\begin{split}
 G_{\{d_i\}}&=B_4(a_j,b,a_k),\\
 G_{\{a_j,a_k\}}&=
 \langle a_1,a_2,a_3,q_i\mid [a_r,a_s]=1,
 q_i a_jq_i^{-1}=a_k,\ q_i a_kq_i^{-1}=a_j,
 (q_i a_i)^2=a_ja_k\rangle,\\
 G_{\{a_j,a_k,d_i\}}&=
 \langle a_j,a_k,q_i\mid[a_j,a_k]=1,
 q_i a_jq_i^{-1}=a_k,\ q_i a_kq_i^{-1}=a_j\rangle,\\
 G_{\{a_1,a_2,a_3\}}&=\mathbb Z^3(a_1,a_2,a_3).
\end{split}                                                       \tag{D.23}
\]
In both middle lines $q_i^2=d_i$; in the third line this
is the separating meridian, not an imposed order-two
relation.  These four families give ten of the twenty-two
vertex-orbit stabilizers.

The complete finite flag index has $22$ vertices, $53$ edges
and $33$ triangles.  Thus the barycentric complex of groups
has $108$ objects, $304$ face arrows and $198$ coherence
triangles.  Every orbit and its actual retained-edge subsets
is listed in \path{code/d_p2_elliptic_three_partial.json}.
The remaining named local maps are still required for the
full dictionary in D.24; these counts do not assert their
completion.
\end{result}
\begin{proof}
Soroko's Proposition~8 and Corollary~9, including their
complete proofs on pages~3--5, identify the boundary-framed
three-holed torus with
\[
 (A(D_4)\times\mathbb Z^3(z_1,z_2,z_3))
       /\langle z_1z_2z_3=\Delta\rangle.
\]
Successive capping kills exactly the three boundary twists
and gives (D.22).  The proof eliminates the other Gervais
generators by the degenerate star relations; in particular
the pair twist opposite $a_i$ becomes $(a_jba_k)^4$.
The figure and those identifications were opened and checked
\cite[Proposition~8 and Corollary~9]{D_Soroko2020}.

Here is the exact algorithm used for subsequent word checks.
Realize the Coxeter group $W(D_4)$ as the $192$ even signed
permutations of four coordinates.  Breadth-first enumeration
supplies the length and a reduced word for each simple element.
The longest element is $-I$, of length $12$, and its positive
lift is the central $\Delta$.  For adjacent simple factors
$x,y$, set
\[
 c=\gcd_{\rm left}(x^{-1}\Delta,y),\qquad
 (x,y)\longmapsto(xc,c^{-1}y).
\]
Each product is length-additive, so this is a positive-monoid
equality.  Repeated transfers terminate: the sequence of
factor lengths strictly increases lexicographically while
its total length is fixed.  Remove identity factors and
central $\Delta$ factors.  The adjacent-gcd criterion is
exactly Dehornoy--Paris Lemma~8.4, whose proof was read;
it identifies the remaining sequence with the unique
greedy normal form \cite[Sections~7--8]{D_DehornoyParis1999}.
A negative atom is replaced by
$\Delta^{-1}(x^{-1}\Delta)$.  Centrality permits collecting
these powers, and cancellativity gives the unique form
$\Delta^m x_1\cdots x_r$ with proper greedy factors.
Two such words agree modulo the central subgroup precisely
when their proper factors agree.

The program computes all $192^2=36864$ left gcds and checks
their uniqueness.  It executes the six Artin relations, four
centrality relations and $400$ seeded contextual rewriting
and normal-form checks.  This supplies a full word algorithm,
not a test in a finite quotient.

Cutting $d_i$ leaves a torus with one marked point and one
boundary, and a pair of pants.  The former has the
four-strand group in (D.23), with boundary twist $d_i$;
this proves the first line, including its meridian.
Cutting $a_j,a_k$ gives a pair of pants and a four-holed
sphere.  The subgroup fixing both curves has kernel
$\mathbb Z^2(a_j,a_k)$ and quotient the pure four-point
group.  The actual twists $a_i,d_i$ are its two standard
intersecting generators.  They commute with the two node
twists and satisfy no further relation, so this subgroup is
\[
             \mathbb Z^2(a_j,a_k)\times F(a_i,d_i).
\]
The only additional graph permutation exchanges the two
node curves.  Its chosen lift is the half twist $q_i$ of
the four-chain subgroup, satisfying $q_i^2=d_i$ and
exchanging $a_j,a_k$.  The star relation gives
\[
                         (q_i a_i)^2=a_ja_k.
\]
The program verifies all these exact words for each of the
three values of $i$; before imposing $\Delta=1$ the last
comparison has precisely one central $\Delta$ factor.
It specifies the conjugation on the remaining free
generator and hence the full index-two extension, proving
the second line of (D.23).

Requiring the additional curve $d_i$ replaces the free
four-point group by its cyclic twist stabilizer.  The
same $q_i$ preserves it, and $q_i^2=d_i$ gives the third
line.  Finally the three complementary pants of
$\{a_1,a_2,a_3\}$ each carry a different fixed label.
Their dual triangle has no nonidentity automorphism
preserving these labels; its stabilizer is exactly its
three independent meridian twists.  This proves the last
line without adding a spurious edge-permutation factor.

For the flag count, use every stable graph of type $(1,3)$
with at least one edge.  For each graph enumerate every
strict chain of nonempty retained-edge subsets, and take
orbits under its full edge-permutation group, realized
by vertex permutations preserving genera and all leg
labels.  Contracting the complementary edges supplies
the actual graph of every face.  This procedure gives
$22,53,33$ orbits in dimensions $0,1,2$.
Each edge has two face arrows; each triangle has six
proper face arrows and six vertex--edge flags.  Hence
there are $2\cdot53+6\cdot33=304$ arrows and
$6\cdot33=198$ coherence triangles.  Parallel-edge
permutations are retained in the enumeration, and
loop half-edge reversals act trivially on these subset
flags.  The saved orbit representatives certify these
counts individually.
\end{proof}
